\documentclass[11pt]{article}
\usepackage{amsmath}
\usepackage[hidelinks]{hyperref}
\title{Which robots need the moving density notice}
\date{27 November 2026}
\begin{document}
\maketitle
\paragraph{Audit.}
Fix the two-block corridor and homogeneous rest-to-rest dynamics
of PR~711. A specified recipient set gets a common identifying
notice at time \(\tau\). Other robots cannot distinguish the
worlds, including indirectly through observing a neighbor's
reaction; this is an explicit observation restriction. The
result characterizes recipient dependence, not radio delivery
or how to enforce that restriction in a physical formation.
Every clearance and modeled link bound holds throughout.
The initial gaps are \(g\), the gap cap is \(R\), and write
\(h=R-g\ge0\) for the slack per gap. Each block has \(m\)
gaps; the terminal demand adds \(\delta\) to whichever
block is identified. Let \(P=B(T)\), \(Q=B(T-\tau)\).

\paragraph{Exact boundary.}
If the middle robot is not informed, the exact maximum
\(\delta\) is \(\min(mh,P)\), even if every other robot
receives the notice. If the middle is informed, let
\(\ell,r\) be the distances in gaps to its closest
uninformed interior neighbor on each side, using \(m\)
on a side with none. The outside boundaries can always be
preprogrammed and need not be counted as recipients. Then
\[
 \delta_{\max}
   =\min\left(mh,\;P+\min\left(Q,\tfrac{h}{2}
                              \min(\ell,r)\right)\right).
\]
In particular, for \(\delta>P\) and \(h>0\), the required
informed neighborhood runs at least through the first
\(q-1\) robots on both sides of the middle, with
\(q=\lceil 2(\delta-P)/h\rceil\), in addition to the middle.
These \(2q-1\) recipients suffice when the other physical
and temporal bounds hold. If \(\delta\le P\), none is needed
if the cap permits the demand.

\paragraph{Lower bound on required information.}
Add the left-world and right-world terminal span
requirements as in PR~711. The two outside displacements
can contribute at most \(2P\), so the half-difference of
the middle's positions must be at least \(\delta-P\).
This half-difference is zero if its local views are shared,
and at most \(Q\) otherwise. Across an uninformed robot
\(i\), the same terminal position must be retained in both
worlds. The span to that robot in either world belongs to
one fixed interval of width \(|i-m|h\), because each of
its gaps belongs to \([g,R]\). Hence the half-difference of
the middle is at most \(|i-m|h/2\). The hard closed
geometric cap also gives \(\delta\le mh\). In the case
there is no uninformed interior neighbor on a side,
the inserted bound \(mh/2\) cannot cut down the maximum:
if \(Q>mh/2\), then \(P\ge Q\) already makes that bound
redundant once \(\delta\le mh\) is enforced.

\paragraph{Construction without informing outsiders.}
If \(\delta\le P\), expand both blocks by \(\delta\) and stay
blind. Otherwise let \(p=P\), \(z=\delta-p\). The outside
boundaries move away from the middle by \(p f(t)\), where
\(f\) is the full-horizon displacement profile normalized to
finish at 1. The middle moves by \(\pm z\,b(t)\), where \(b\)
is the similarly normalized profile over the shorter horizon,
zero before notice. Since \(0<z\le Q\le p\), these scaled
profiles respect the original actuator bounds. On a side
with nearest blind index \(\ell\), choose preparation
\[
            s\in[z,\ell h-z]\cap[p-(m-\ell)h,p] .
\]
The inequalities \(2z\le\ell h\), \(p+z\le mh\), and
\(z\le p\) make this interval nonempty. The blind robot
at that index moves outward by \(s f(t)\) in both worlds.
Interpolate the displacement linearly among the remaining
robots between it and the outer boundary, and among the
informed robots between it and the middle. If \(\ell=m\),
the two exterior positions coincide by the imposed choice
\(s=p\). Treat the right side in the same manner using \(r\).
The near-span extra is \(s f(t)\pm z b(t)\).
It is in \([0,\ell h]\) because \(0\le b(t)\le f(t)\le1\);
the far span similarly fits its cap, and uniform
interpolation gives the same result for every individual
gap. Interior derivatives are convex combinations of
the already bounded boundary derivatives. The active
terminal block gains exactly \(p+z=\delta\).

\paragraph{Order of normalized profiles.}
For \(0<t<T\), let \(u(t)=\min(at,a(T-t),v)\).
The shorter envelope's speed is
\(w(t)=\min(u(t),a(t-\tau)_+)\).
Thus \(w(t)/u(t)\) is nondecreasing: before \(\tau\) it is
zero, and afterward it is the minimum of 1 and the maximum
of \((t-\tau)/t,\ (t-\tau)/(T-t),\ a(t-\tau)/v\).
All three ratios are nondecreasing on that interval.
For any time split, integrate \(w(x)u(y)\le u(x)w(y)\)
 over \(x\le t\le y\) in the horizon. Dividing by the
positive integrals proves that the accumulated normalized
short profile is no larger than the full one. When the
short phase has no travel allowance, the moving branch
is unnecessary. This supplies the continuous hard-gap
bound, rather than extrapolating sampled gaps.

\paragraph{Checks and limits.}
With \(m=2,h=3/2,a=2,T=2,\tau=1/2\) and nonbinding speed,
the full geometric demand 3 can be met by the middle
and its two immediate neighbors. Informing only the
middle reduces the exact maximum to \(11/4\); informing
others but not the middle reduces it to 2. Exact terminal
difference-feasibility checks agree on 73,662 inputs, and
12,008 dyadic samples check the normalized profiles.
Four Lean lemmas cover the key real interval steps.
The strongest direct planner with the same observation
restriction ties. A robot that can sense the middle's
branch in time would not be in the blind class used here.
Acquiring and propagating that information, later or
inconsistent delivery, and any independent formation
evaluation remain untested.
\end{document}
